\documentclass[11pt]{article}
\usepackage[margin=0.82in]{geometry}
\usepackage{amsmath,amssymb,amsthm,graphicx,microtype}
\usepackage[hidelinks]{hyperref}

\title{Exact polynomial content of an atomic singular-inner\\
borderline de Branges--Rovnyak family}
\author{Partial result for arXiv:2405.07016}
\date{13 August 2026}

\newtheorem{theorem}{Theorem}
\newtheorem{lemma}{Lemma}
\newcommand{\D}{\mathbb D}
\newcommand{\C}{\mathbb C}
\newcommand{\Hh}{\mathcal H}

\begin{document}
\maketitle

\begin{center}
\fbox{\parbox{0.92\textwidth}{\textbf{Status: substantial partial result, likely valid.}
For a concrete infinite-dimensional family of extreme singular-inner symbols,
the polynomials in the borderline Schur-power space are classified exactly.
The source's existential question for arbitrary extreme symbols remains open.}}
\end{center}

\section*{1. Source question}

Aleman and Dahlin study generalized de Branges--Rovnyak spaces and, at the
end of arXiv:2405.07016v3, ask the following question \cite{AlemanDahlin}.
Here $\beta=m$ and $d=1$, so the relevant kernel is
\begin{equation}\label{eq:targetkernel}
 K_{b,m}(z,w)=
 \left(\frac{1-b(z)\overline{b(w)}}{1-z\bar w}\right)^m.
\end{equation}

\begin{figure}[h]
\centering
\includegraphics[width=0.96\textwidth]{figures/open_problem_crop.png}
\caption{Question 1 on page 25 of arXiv:2405.07016v3.}
\end{figure}

In transcription: does there exist an integer $m>1$ and an extreme point $b$
of the unit ball of $H^\infty$ such that $\Hh(K_{b,m})$ contains every
polynomial, and, if so, are the polynomials dense?  The source proves the
positive conclusion for non-extreme $b$, but the displayed example $b(z)=z$
gives only the constants at the borderline.

\section*{2. Result and proof intuition}

For $a>0$, let
\begin{equation}\label{eq:symbol}
 S_a(z)=\exp\!\left(-a\frac{1+z}{1-z}\right),
 \qquad b_{N,a}(z)=z^N S_a(z),\qquad N\geq1.
\end{equation}
The function $S_a$ is the singular inner function associated with one atom at
$1$, so $b_{N,a}$ is inner and hence extreme.

\paragraph{Proof intuition.}
The model kernel for $z^NS_a$ splits into a finite polynomial block and a
shifted copy of the kernel for $S_a$.  Expanding its $m$-th power shows at once
that all monomials through degree $m(N-1)$ occur.  To show that no later
monomial can be hidden in the Hilbert completion, move to the right half-plane.
There $S_a$ becomes $e^{-as}$, and every element of the $j$-th kernel power is
$(s+1)^j$ times the Laplace transform of a function supported in $[0,ja]$.
That Laplace transform is entire.  Applied to a putative polynomial, this
forces its pole at the Cayley point $s=-1$ to have bounded order, giving the
same sharp degree cutoff.  This settles an infinite-dimensional extreme
family but does not rule out more irregular extreme non-inner symbols.

\begin{theorem}\label{thm:main}
For $a>0$ and integers $N,m\geq1$,
\[
 \Hh(K_{b_{N,a},m})\cap\C[z]
 =\{p\in\C[z]:\deg p\leq m(N-1)\}.
\]
In particular, for every $m>1$ the infinite-dimensional space
$\Hh(K_{zS_a,m})$ contains no nonconstant polynomial.  In all these spaces the
polynomials fail to be dense.
\end{theorem}

\begin{proof}
Write $t=z\bar w$ and $P_N(t)=1+t+\cdots+t^{N-1}$.  The elementary product
identity for model kernels gives
\begin{equation}\label{eq:split}
 K_{z^NS_a}(z,w)=P_N(t)+t^N K_{S_a}(z,w),
\end{equation}
where $K_u(z,w)=(1-u(z)\overline{u(w)})/(1-z\bar w)$.  Consequently,
\begin{equation}\label{eq:powerexpand}
 K_{z^NS_a}(z,w)^m
 =\sum_{j=0}^m\binom mj P_N(t)^{m-j}t^{Nj}K_{S_a}(z,w)^j.
\end{equation}
The $j=0$ term is a diagonal polynomial kernel whose coefficient is strictly
positive in every degree from $0$ to $m(N-1)$.  Thus every polynomial of
degree at most $m(N-1)$ belongs to the target space.

We next describe the other summands without an algebraic-product closure
assumption.  Use the Cayley coordinate
\[
 z=\phi(s)=\frac{s-1}{s+1},\qquad \operatorname{Re}s>0.
\]
Since $S_a(\phi(s))=e^{-as}$, direct calculation gives
\begin{equation}\label{eq:halfkernel}
 K_{S_a}(\phi(s),\phi(t))
 =\frac{(s+1)(\bar t+1)}2
   \frac{1-e^{-a(s+\bar t)}}{s+\bar t}.
\end{equation}
For $j\geq1$,
\begin{equation}\label{eq:laplacekernel}
 \left(\frac{1-e^{-a(s+\bar t)}}{s+\bar t}\right)^j
 =\int_0^{ja}e^{-u(s+\bar t)}w_j(u)\,du,
 \qquad w_j=\mathbf1_{[0,a]}^{*j}.
\end{equation}
The feature-map construction of the RKHS in \eqref{eq:laplacekernel} shows
that every $f\in\Hh(K_{S_a}^j)$ has the form
\begin{equation}\label{eq:featureform}
 f(\phi(s))=(s+1)^jF(s),
\end{equation}
up to a fixed nonzero scalar, where
\[
 F(s)=\int_0^{ja}h(u)e^{-su}w_j(u)\,du
 \quad\text{for some }h\in L^2(w_j(u)\,du).
\]
In particular $F$ is entire: compact support and Cauchy--Schwarz give locally
uniform convergence of the integral on all of $\C$.  This feature-space
description covers the completed RKHS, not only finite products.

Expand $P_N^{m-j}$ in \eqref{eq:powerexpand}.  The standard finite-sum rule
for RKHSs \cite{Aronszajn} implies that any polynomial
$p\in\Hh(K_{b_{N,a},m})$ admits a decomposition
\begin{equation}\label{eq:decomp}
 p=p_0+\sum_{j=1}^m\sum_{\ell=0}^{(m-j)(N-1)}
 z^{Nj+\ell}f_{j,\ell},
\end{equation}
where $\deg p_0\leq m(N-1)$ and
$f_{j,\ell}\in\Hh(K_{S_a}^j)$.  Put $q=p-p_0$.  Every term on the right is
divisible by $z^N$, so $q=z^Nh$ for a polynomial $h$.  Substituting
\eqref{eq:featureform} into \eqref{eq:decomp} and dividing by $z^N$ yields
\begin{equation}\label{eq:pole}
 h(\phi(s))=
 \sum_{j,\ell}(s-1)^{r_{j,\ell}}
 (s+1)^{j-r_{j,\ell}}F_{j,\ell}(s),
 \qquad r_{j,\ell}=N(j-1)+\ell.
\end{equation}
Set $R=m(N-1)-N$.  For every index in \eqref{eq:pole},
\[
 r_{j,\ell}-j=(N-1)j-N+\ell
 \leq m(N-1)-N=R.
\]
It follows that $(s+1)^Rh(\phi(s))$ extends to an entire function.  If the
polynomial $h$ is nonzero of degree $d$, then
$h((s-1)/(s+1))$ has a pole of exact order $d$ at $s=-1$.  Hence $R<0$
forces $h=0$, whereas $R\geq0$ forces $d\leq R$.  In either case,
\[
 \deg q=N+\deg h\leq N+R=m(N-1)
\]
whenever $q\ne0$.  The same degree bound already holds for $p_0$, proving the
claimed equality.

Finally, $S_a$ is not a finite Blaschke product, so its model space is
infinite-dimensional; \eqref{eq:split} and \eqref{eq:powerexpand} then make
$\Hh(K_{b_{N,a},m})$ infinite-dimensional.  Its polynomial subspace is finite
dimensional by the proved classification and therefore cannot be dense.
\end{proof}

\section*{3. Verification, scope, and novelty}

The proof was audited at three points: the positive-kernel expansion
\eqref{eq:powerexpand}, the compactly supported feature representation
\eqref{eq:laplacekernel}, and the maximal exponent
$r_{j,\ell}-j=R$.  The last identity was also checked mechanically over a
grid of integer parameters by \texttt{code/exponent\_audit.py}.  Two separate
finite-section programs tested an infinite Blaschke candidate and an explicit
non-inner extreme outer candidate; both showed rapidly growing lower bounds
for missing monomial norms.  Those numerical diagnostics are included in
\texttt{code/}, but they are not used in Theorem~\ref{thm:main}.

This is a partial result only.  It gives exact negative answers for a natural
atomic singular-inner family, including infinite-dimensional examples, but
does not decide whether some other extreme $b$ makes all polynomials belong.
The strongest attempted extension to every inner symbol produces an
antiunitary conjugation on finite product sums, but identifying its extension
with the formal boundary conjugate is blocked by the lack of boundary control
from the ambient weighted Bergman norm.

The run indexes contained no duplicate.  Bounded searches through 13 August
2026 covered the exact question text, the paper title and authors, current
arXiv versions, the 2025 JFA publication, the first author's 2026 thesis, and
the higher-order de Branges--Rovnyak literature.  No resolution of Question 1
or this atomic singular-inner classification was found.  Novelty outside the
run is not asserted pending expert review.

\paragraph{Human-review recommendation.}
Check the RKHS finite-sum decomposition in \eqref{eq:decomp}, the feature-space
surjectivity used in \eqref{eq:featureform}, and the pole-order argument when
$R=-1$.  Preserve the partial scope: the result supplies a sharp negative
family, not a solution of the existential question.

\begin{thebibliography}{9}
\bibitem{AlemanDahlin}
A. Aleman and F. Dahlin,
\emph{Generalized de Branges--Rovnyak spaces},
J. Funct. Anal. 288 (2025), no. 11, 110860;
arXiv:2405.07016.

\bibitem{Aronszajn}
N. Aronszajn,
\emph{Theory of reproducing kernels},
Trans. Amer. Math. Soc. 68 (1950), 337--404.
\end{thebibliography}

\end{document}
